\chapter{对偶控制与热信息探针}
\label{ch:ati}

一次边界变化同时改两个量：参数后验和品质泛函。二者可以驱动同一阀门。把每一次喷蒸汽都称为认知，或在方差没有下降时仍提前停机，都会把交叉表读错。探针因此写成有限时长的执行器轨迹，已经运行的时间表是 \([150,190)\) 秒的两段脉冲；四种策略在同一加热律上给出不同的停机几何。

\section{探针的定义}

\begin{definition}[热信息探针]
探针是一段有限时长 \(\tau\le\tau_{\max}\) 的执行器轨迹 \(a\in\cA\)，满足三条：设计目标是增大\eqref{eq:fisher}的 \(\log\det\)，或等价地减小决策量的后验熵；对表面褐变、质量损失和热力进度的预测损伤低于预算 \(\eta\)；结束后执行器回到探针前的名义轨迹，除非安全层已经改写该轨迹。
\end{definition}

不满足第一条的功率变化是加工动作。高斯近似下，决策状态的微分熵由协方差的行列式决定。对候选探针 \(a\)，
\begin{equation}
\label{eq:ati}
G(a)=\frac12\log\det\vect{\Sigma}_{z,k}
-\frac12\log\det\vect{\Sigma}_{z,k+\tau}(a).
\end{equation}
两种信息形式都没有在线计算。对偶策略用时间窗代替定义：\([150,190)\) 秒且方差大于 0.25 时进入探针模式。前 20 s 蒸汽全开、排汽关小，后 20 s 风机升高到 1.3。这是一条信息动作的时间表：顺序保留了先冷凝、后对流，时长固定为 40 s，打分用热阻是否更接近参照来代替 \(\log\det\)。

窗结束时，若新的热阻比旧值更接近由真值温差构成的参照，方差乘 \(0.35\)。参照
\[
r_{\mathrm{proxy}}=\mathrm{clip}(\Delta T/8+1.2+2.5 S,\,0.8,\,8)
\]
用了表面与冷点的真值温差，实施到机器上时改成可安装通道。方差没有下降，停止时刻就留在动作之前。厚度超过 \(40\,\mathrm{mm}\) 且方差仍高于 \(0.45\) 时，停机判据要求方差先下降，点估计单独越过阈值仍保持加热。

\section{两段脉冲的秩}

\begin{proposition}[第二输入方向]
\label{prop:pulse}
线性化后，仅 \(\delta T_\infty\) 时输入矩阵为单列。在\eqref{eq:twonode}里 \(q_{\mathrm{lat}}\) 与 \(T_\infty\) 都只进入第一式，两列输入共线，变号不增加输入矩阵的秩。可分开的是参数列：离开平衡之后，只改变 \(R_{\mathrm{ext}}\) 的气流段使 \((R_{\mathrm{ext}},R_{\mathrm{int}})\) 的灵敏度列线性无关。
\end{proposition}

\begin{proof}
\(T_\infty\) 的输入列是 \((1/(C_1 R_{\mathrm{ext}}),\,0)\)，\(q_{\mathrm{lat}}\) 的输入列是 \((1/C_1,\,0)\)。第二分量都是 0，潜热变号只改变倍数的符号，秩仍是 1。气流段把 \(R_{\mathrm{ext}}\) 改为另一个正数，对 \(R_{\mathrm{ext}}\) 的偏导数含有 \(-(T_\infty-T_1)/R_{\mathrm{ext}}^2\)，对 \(R_{\mathrm{int}}\) 的偏导数含有两节点温差。冷凝段对 \(R_{\mathrm{ext}}\) 的偏导数不包含潜热本身。因此气流段的灵敏度列含有一个冷凝段列空间里没有的分量。
\end{proof}

断言在两节点上比较两段轨迹与单独蒸汽的内部滞后，得到 40.846 与 38.912。判据改为差值大于 1 之后通过。相对倍数 \(1.050\) 小于 1.15，旧的相对判据没有通过。通过只说明第二段在该组常数上改变了滞后。

\section{探针窗口内的冷点上界}

\begin{proposition}[冷点增量]
\label{prop:budget}
记 \(\tau_2=R_{\mathrm{int}}C_2\)。当 \(\tau\le\tau_2/10\) 且表面温度增量不超过 \(\Delta T_s\) 时，
\[
|\Delta T_2|\le\frac{\tau}{\tau_2}\Delta T_s,\qquad
|\Delta F|\le\tau\,10^{(T_2(0)+\Delta T_2-T_{\mathrm{ref}})/z}.
\]
\end{proposition}

\begin{proof}
第二节点是线性滤波器 \(\dot T_2=(T_1-T_2)/\tau_2\)。在窗口 \([0,\tau]\) 上把 \(T_1-T_2\) 放大到不超过 \(\Delta T_s\)，积分得所述上界。前提 \(\tau\le\tau_2/10\) 保证温差还没有被冷点的上升吃掉。
\end{proof}

损伤预算因此写在表面褐变和表面含水上。用冷点温度否决探针，会放过表面已经越界的情况。断言取 \(\tau_2=30\)、\(\tau=3\)、\(\Delta T_s=12\)，上界 1.2，离散步进得到 1.143。这组常数属于命题核对，S1 的窗口是 40 s。种子 7 在 40 s 里冷点从 \(8.01^\circ\mathrm{C}\) 到 \(8.05^\circ\mathrm{C}\)，表面从 \(89.4^\circ\mathrm{C}\) 到 \(100.3^\circ\mathrm{C}\)，短窗口把热量留在近表面。

40 s 探针窗在 30 mm 上的 Fourier 数为 \(8.9\times 10^{-3}\)，远小于 1。Biot 数在风机 0.4 时为 1.5，探针段风机 1.3 时升到 3.7，表面梯度更陡。这是两段探针要把外部热阻和内部热阻分开的尺度理由。图\ref{fig:fobi6}把 \(\mathrm{Fo}(t)\) 与 \(\mathrm{Bi}(u_{\mathrm{fan}})\) 画在核常数下；完整讨论见第\ref{ch:section}章。

\begin{figure}[htbp]
\centering
\includegraphics[width=0.82\textwidth]{X_fo_bi.png}
\caption{核常数 \(\alpha=2.0\times 10^{-7}\,\mathrm{m}^2/\mathrm{s}\)、\(k=0.68\,\mathrm{W}/(\mathrm{m\cdot K})\)、\(h=12+55u_{\mathrm{fan}}\) 下的 Fourier 数与 Biot 数。}
\label{fig:fobi6}
\end{figure}

\section{四种策略}

加热请求
\begin{equation}
\label{eq:ureq}
u_{\mathrm{req}}=\mathrm{clip}\big(0.35+(75+\mu-\hat T_c)/30,\,0.15,\,1\big),
\end{equation}
裕量 \(\mu=6+10\sqrt{\mathrm{Var}}\,(L/0.03)\)。式子只看见估计温度、方差和厚度。停机要求 \(\hat T_c>75+\mu\) 且时刻已经超过 240 s。

\begin{enumerate}
\item \textbf{开环。} 按厚度给菜单时间，中途不改。15 mm、30 mm、50 mm 分别是 900 s、1170 s、1530 s。
\item \textbf{腔温比例。} \(u=\mathrm{clip}(0.04(160-y_c)+0.6,0,1)\)。只跟踪带噪腔温，不停机。
\item \textbf{无探针。} 使用\eqref{eq:ureq}，没有蒸汽窗。
\item \textbf{对偶。} 在无探针加热律上叠加固定探针窗，厚件上另有方差否决。
\end{enumerate}

无探针和对偶共用同一条饱和线性加热律，差别在 40 s 窗口和其后的热阻更新。交叉表里无探针的提前比例 \(0.944\)、对偶 \(0.902\)：估计温度经常先于真值进度越过阈值，提前来自跟踪器把点估计送到停机门槛之前。开环把停止时刻写成厚度的函数，15/30/50 mm 分别钉在 900/1170/1530 s；腔温比例把停止时刻写成时域终点 3300 s。四条策略因此在同一积分器上给出四种停机几何：时钟、腔温、点估计、点估计加短窗。

加热请求\eqref{eq:ureq} 的裕量随方差和厚度放大。方差仍为先验 0.64、厚度 30 mm 时 \(\mu\approx 14\,\mathrm{K}\)，停机门槛大约 \(89^\circ\mathrm{C}\)；方差落到 0.224 之后 \(\mu\) 降到约 \(11\,\mathrm{K}\)。厚件上同一公式把门槛再抬高，点估计更难单独宣布停机。这是对偶策略在 50 mm 上提前比例为 1、中位仍为负的跟踪器侧原因：门槛跟着厚度走，导热时间尺度却走不完。

图\ref{fig:seed7}把一条轨迹读完整。蒸汽峰值只出现在探针窗。1680 s 仍是品质加热，请求 0.476，真值冷点 \(70.21^\circ\mathrm{C}\)，\(\alpha_m=0.558\)。品质段结束于 1864 s，当时冷点 \(77.70^\circ\mathrm{C}\)，\(\alpha_m=0.995\)，\(e_{\alpha}=+0.145\)。保持从 1872 s 开始。评价用停机时刻的 \(\alpha_m\)：终点进度已经饱和到 1，剖面已经拉平，那是保持段的结果。完整抽样见附录\ref{app:trace}。

\begin{figure}[!htbp]
\centering
\includegraphics[width=0.82\textwidth]{P_seed7_panels.png}
\caption{种子 7、肉糜 30 mm、对偶策略。阴影为探针窗口。方差在窗后由 0.64 降到 0.224；停机时刻进度为 0.995，保持段把剖面拉平。}
\label{fig:seed7}
\end{figure}

图\ref{fig:pdt}把表面减冷点、估计减冷点画在一起。表面温差在品质段升到 \(120\,\mathrm{K}\) 量级；估计减冷点在中段为正，切入保持后落到 \(-20\,\mathrm{K}\)。估计可以高于冷点，给出的是带滞后的点跟踪，下置信界需要后验分位。

探针窗是一条可以在日志里核对的时间表。种子 7 在 150 s 进入探针，蒸汽从 0.15 拉到 1.0，排汽关小；170 s 风机升到 1.3；190 s 窗口结束，蒸汽回到 0.15。40 s 内冷点从 \(8.01^\circ\mathrm{C}\) 升到 \(8.05^\circ\mathrm{C}\)，表面从 \(89.4^\circ\mathrm{C}\) 升到 \(100.3^\circ\mathrm{C}\)，热量留在近表面。窗后若新热阻更接近参照，方差从 0.64 乘到 0.224。这一乘子在 15、30、50 mm 上相同：探针买到的是方差，这块食品在 55 min 内能否到达设计线，还要看扩散时间。

信息动作与品质动作因此在时间轴上分开。150 s 之前四个执行量按品质加热走，加热请求由估计温度、方差和厚度决定；150--190 s 蒸汽和风机按时间表离开平衡，加热请求仍由同一条饱和律给出，只是边界被探针改写；190 s 之后蒸汽回到基线 0.15，风机回到 0.4，方差若已下降则停机门槛跟着下降。日志里要能查出这一时刻是哪一个函数发出的指令：探针窗内蒸汽列跳到 1.0，窗外蒸汽列回到 0.15，加热列始终由\eqref{eq:ureq} 给出。可分离落实为这条时间表。在线最大化 \(\log\det\) 的求解器留在规格里。

\begin{figure}[htbp]
\centering
\includegraphics[width=0.82\textwidth]{P_dt.png}
\caption{种子 7：表面减冷点、估计减冷点。估计在中段高于冷点。}
\label{fig:pdt}
\end{figure}
\FloatBarrier
